\documentclass[11pt]{article}
\usepackage{fontspec}
\setmainfont{Times New Roman}
\setsansfont{Arial}
\setmonofont{Consolas}
\usepackage{amsmath,amssymb,mathtools}
\usepackage{geometry}
\usepackage{microtype}
\geometry{a4paper,margin=2.15cm}
\setlength{\parindent}{1.15em}
\setlength{\parskip}{0pt}
\emergencystretch=180em
\allowdisplaybreaks
\sloppy
\hbadness=10000
\vbadness=10000
\hfuzz=3pt
\vfuzz=10pt
\raggedbottom
\begin{document}
% Unit: paper35_source_fidelity_v001.
% Language lane: Interslavic Cyrillic.
% Direct source: sources/paper35/source_fidelity/Noether_Paper35_ORIGINAL_MathNet600_R124plusP40_repaired_witness_v001.tex.
\section*{35. О максималных областјах целочислных функциј}
\emph{Рец. Соц. Матх. Мосцоу 36 (1929), с. 65--72. Изворно-верны превод по Р124+П40/МатхНет600 поправеном немецком сведку.}

Следујуче изследованьа близко дотыкајут Хилбертову проблему релативно целых функциј\footnote{Д. Хилберт, Матхематисцхе Проблеме (Гöтт. Нацхр. 1900, с. 253--297), проблем 14.}; оне изходет из случајного впрашаньа Е. Хецке, кторы мог непосредно рачунско ријешити специалны случај, потребны јему како помочна лема\footnote{Е. Хецке, Üбер дие Конструктион релатив-Абелсцхер Захлкöрпер дурцх Модулфунктионен вон зwеи Вариабелн (Матх. Анн. 74 (1913), с. 465--510), §2. У Хецке конечна интегритетна област \(\mathfrak{I}\) породжена јест трома квоциентами степеновых редов в двух променных, меджу кторыми јест алгебраична релација.}.

Впрашанье јест следујуче. Под \textit{максималноју областју} целочислных функциј в \(x_1,\ldots,x_n\) --- полиномов, степеновых редов или квоциентов такых --- внутри даној области \(\mathfrak{Z}\) разумеју таку област, коју уже не можно разширити додаваньем целочислных именовательев: ако \(a\cdot g(x)\) принадлежи јеј, кде \(g(x)\in\mathfrak{Z}\), тогда всегда принадлежи и \(g(x)\); ту \(a\) значи ненулево цело рационално число. Всака интегритетна област \(\mathfrak{I}\) в \(\mathfrak{Z}\) једнозначно пораджа максималну област \(M(\mathfrak{I})\), која состојí из всих функциј \(g(x)\) из \(\mathfrak{Z}\), за кторе јест најменеј једно \(a\ne0\) тако, же \(a\cdot g(x)\in\mathfrak{I}\).

Идé о впрашанье, что можно ректи о именовательах \(a\) тој максималној области \(M(\mathfrak{I})\), ако почетна интегритетна област \(\mathfrak{I}\) была конечна, то јест состојала из всих целочислных полиномов конечно многих функциј \(f_1(x),\ldots,f_r(x)\) из \(\mathfrak{Z}\). При том за \(\mathfrak{Z}\) треба взяти все целочислне полиномы или степенове реды в \(x_1,\ldots,x_n\), респективно таке квоциенты, кторе не имајут инех именовательев кроме тых, кторе се јавјајут в \(f(x)\), и јих степенов.

Показујем, же под додатным условјем в случају степеновых редов или јих квоциентов --- именно, же меджу \(f(x)\) сут само алгебраичне зависимости --- числа \(a\) можно всегда избирати тако, же в них вцходит само конечно много различных првочисел, иако обчно в всаких степеньах.

Последнье јест непосредно јасно: ако \(a\cdot g(x)\in\mathfrak{I}\), але \(a^\varrho g(x)\) не принадлежи за никакы експонент \(\varrho\), тогда појавјајут се все степены \(a^\varrho\); например при \(\mathfrak{I}=[2x]\) имајемо \(x=(2x)/2\), \(x^2=(2x)^2/2^2\).

Впрашанье ограничености можно формуловати само тако: можно ли дати, обчно бесконечну, систему породительев за \(M(\mathfrak{I})\), тако же в соответных именовательах експоненты првочисел оставајут ограничене? Понеже идé само о конечно много првочислах, то впрашанье по познаных доказах\footnote{Срав. Хилберт, тамже, или подробнејше изложенье у Е. Ноетхер, Дие Ендлицхкеит дес Сыстемс дер ганззахлиген Инвариантен бинäрер Формен (Гöтт. Нацхр. 1919, с. 138--156), §5.} јест исто с впрашаньом конечности \(M(\mathfrak{I})\), то јест с Хилбертово проблемо релативно целых функциј в целочислној форме. То впрашанье с ту употребјеными методами не могу решити\footnote{Тако, например, в случају целочислных пројективных инвариантов основној системы форм показано јест, же оны се могу представити чрез полну систему инвариантов в обычном смислу, с само конечно много различными првочислами в именовательу. То не следује из обычној методы \(\Omega\)-процеса. Але ничто не јест речено о самој целочислној конечности, ктора доселе познана јест само за бинарне основне формы. Срав. ноту цитовану в приметце 3.}.

Доказ, же в именовательах \(a\) вцходит само конечно много различных првочисел \(p\), опира се на том, же всакому такому \(p\) одговарја најменеј једна релација модуло \(p\) меджу функцијами \(f(x)\), кторе пораджајут \(\mathfrak{I}\), релација с целочислными коефициентами, ктора не јест идентично изполньена в \(x\). Другыми словами: ако \(\mathfrak{o}\) значи област целочислных полиномов в неизвестных \(u_1,\ldots,u_r\), тогда \(\mathfrak{I}\) чрез присвојенье \(u_i\mapsto f_i(x)\) јест хомоморфны образ \(\mathfrak{o}\); тут хомоморфизм посредованы јест простым идеалом \(\mathfrak{a}\) из \(\mathfrak{o}\), то јест \(\mathfrak{I}\) изоморфна јест фактор-колцу \(\mathfrak{o}/\mathfrak{a}\). Соответне \(\mathfrak{I}_p\) јест хомоморфны образ \(\mathfrak{o}_p\), кде индекс значи преход к остатковым класам модуло \(p\), можны с изкльученьем конечно многих првочисел в именовательах \(\mathfrak{I}\). Тут хомоморфизм знову посредованы јест простым идеалом \(\mathfrak{c}_p\) из \(\mathfrak{o}_p\), кторы обхвата \(\mathfrak{a}_p\). Тогда и само тогда \(p\) вцходит в именователь \(M(\mathfrak{I})\), когда \(\mathfrak{c}_p\) стаје стриктны надидеал \(\mathfrak{a}_p\). То може настати само ако или трансценденцијны степен \(\mathfrak{I}_p\) се понижа проти \(\mathfrak{I}\), или \(\mathfrak{a}_p\) губи својство простого идеала\footnote{Под "трансценденцијным степеньем" системы разумеје се познано инвариантно число алгебраично независимых функциј, од кторых система алгебраично зависи; система из алгебраично независимых функциј именује се иредуцибилна. Под трансценденцијным степеньем, или димензијеју, простого идеала разумеје се степен јего остаткового польа. Просты идеал не има стриктного простого делительа истого трансценденцијного степена; срав. Б. Л. ван дер Wаерден, Зур Нуллстеллентхеорие дер Полыномидеале, Матх. Анн. 96 (1926), с. 183--208.}. Же то настаје само за конечно много \(p\), следује из того, же и модуло \(p\) алгебраична зависимост влече изчезньенье функционалного детерминанта (§1), и же \(\mathfrak{a}\) јест абсолутны просты идеал и зато губи простоту модуло само конечно многих \(p\)\footnote{Е. Ноетхер, Елиминатионстхеорие унд аллгемеине Идеалтхеорие, Матх. Анн. 90 (1923), с. 229--261, §7. Тврдженье следује чрез елиминацијону теорију из простејшего тврдженьа, же абсолутно иредуцибилны полином губи то својство модуло само конечно многих првочисел; А. Остроwски, Зур аритхметисцхен Тхеорие дер алгебраисцхен Грößен, Гöтт. Нацхр. 1919, с. 279--298, там с. 296. Простејши доказ у Е. Ноетхер, Еин алгебраисцхес Критериум фüр абсолуте Ирредузибилитäт, Матх. Анн. 85 (1922), с. 26--33, §7. Ако \(\mathfrak{I}\) има интегритетну базу, ктора има точно једну функцију вечше неж јеј трансценденцијны степен, како в случају Хецке из приметкы 2, тогда \(\mathfrak{a}\) стаје главны идеал и достачује простејше тврдженье без елиминацијној теорије. Наместо првочисел всуде могу стајати просты идеалы конечного алгебраичного числового польа. Тврдженье јавно остаје валидно и когда \(\mathfrak{a}\) јест нуловы идеал.}.

Треба јеште приметити, же все разсудживаньа с малыми модификацијами оставајут валидне, ако наместо целых рационалных чисел взяти целе числа конечного алгебраичного числового польа и јего просте идеалы (§4).

\bigskip
\noindent\textbf{§ 1. Функционалне детерминанты над областами коефициентов карактеристики \(p\).}

\smallskip
\noindent\textbf{1.}
Нехај \(f_1(x_1,\ldots,x_n),\ldots,f_t(x_1,\ldots,x_n)\) сут рационалне функције или, обчнеје, формално дефиноване степенове реды или квоциенты такых, с коефициентами из польа \(P\) карактеристики \(p\). Изводы \(\partial f_i/\partial x_k\) такоже нехај будут формално дефиноване, при чем все обычне рачунске законы оставајут. Под \(\Omega\) разумеје се алгебраично затворено разширенье польа \(P\).

Тогда, како и при карактеристики нула, вельа:

\noindent\textbf{Теорем.}
\textit{Ако меджу \(f(x)\) јест алгебраична зависимост с коефициентами из \(\Omega\), тогда ранг функционалној матрице \((\partial f_i/\partial x_k)\), \(i=1,\ldots,t\), \(k=1,\ldots,n\), јест менши неж \(t\)}\footnote{В изворној формулацији предпоставјала јесм, же коефициентна област \(P\) јест перфектно полье. Разширенье на либокакве, такоже неперфектне польа, чрез преход од \(P\) к \(\Omega\), должна јесм Б. Л. ван дер Wаерден. Же обратно тврдженье при карактеристики \(p\) не вельа и при перфектном \(P\), показује примјер \(f=x^p\), \(\partial f/\partial x=0\).}.

\noindent\textit{Доказ.}
В полиномној области \(\Omega[u_1,\ldots,u_t]\) нехај \(\mathfrak{m}\) значи просты идеал всих полиномов \(G(u)\), за кторе \(G(f)=0\) идентично в \(x\). По предпоставјеньу \(\mathfrak{m}\) не јест нуловы идеал; нехај \(G(u)\ne0\), зато и \(G(u)\ne\mathrm{const.}\), јест полином најнижшего степена из \(\mathfrak{m}\). Из того особно следује, же \(G(u)\) јест иредуцибилны. Тогда како обычно добивамо
\[
\frac{\partial G}{\partial u_1}\cdot\frac{\partial f_1}{\partial x_k}
+\cdots+
\frac{\partial G}{\partial u_t}\cdot\frac{\partial f_t}{\partial x_k}=0
\quad(k=1,\ldots,n),
\]
одкуд следује, же или ранг функционалној матрице јест менши неж \(t\), или все \(\partial G/\partial u_i\) изчезајут. Але понеже \(G(u)\) был избран како полином најнижшего степена, из того следује изчезньенье всих \(\partial G/\partial u_i\).

Тако имамо \(G(u)=H(u^p,\ldots,u^p)\); а понеже коефициентна област \(\Omega\) јест алгебраично затворено, зато перфектно, полье, далье \(G(u)=H(u)^p=(H(u))^p\), противно избору \(G(u)\). Доказ наравно вельа и за карактеристику нула, кде последньи крок одпада.

\medskip\noindent\textbf{2. Последок.}
Нехај \(F_1(x_1,\ldots,x_n),\ldots,F_t(x_1,\ldots,x_n)\) сут целочислне функције, то јест полиномы, степенове реды или квоциенты такых, с функционалноју матрицеју точно ранга \(t\). Првочисло \(p\) нехај буде избрано тако, же оно не вцходит ни в именователье \(F(x)\), ни в все \(t\)-редне детерминанты функционалној матрице.

Тогда \(F_1(x),\ldots,F_t(x)\) оставајут алгебраично независиме и модуло \(p\).

\noindent\textit{Доказ.}
Нехај \(\bar P\) значи остатково полье по \(p\), а \(f_1(x),\ldots,f_t(x)\) функције из \(\bar P\), кторе настајут из \(F_1(x),\ldots,F_t(x)\), когда целочислне коефициенты заменьају се својими остатковыми класами по \(p\). Функције \(f_i(x)\) сут дефиноване, бо \(p\) не вцходило в именователь \(F_i(x)\). Јих функционална матрица \((\partial f_i/\partial x_k)\) настаје из \((\partial F_i/\partial x_k)\) такоже заменоју целочислных коефициентов остатковыми класами по \(p\); в \((\partial F_i/\partial x_k)\) не појавјајут се ине именователье неж в \(F(x)\). Заради избора \(p\) матрица \((\partial f_i/\partial x_k)\) зацховава ранг \(t\); функције \(f(x)\) сут алгебраично независиме в \(\Omega\), и тем боль в \(\bar P\).

\bigskip
\noindent\textbf{§ 2. Формулација теорема.}

\smallskip
\noindent\textbf{1.}
Како в уводу, нехај \(\mathfrak{I}\) значи интегритетну област, то јест колцо без нуловых делительев, из целочислных функциј в \(x_1,\ldots,x_n\): полиномов или степеновых редов с целыми рационалными коефициентами или квоциентов такых. В случају степеновых редов може се говорити о формално дефинованых или о сходных в некој области; употребльа се само то, же појавјајуче се функције творет интегритетну област.

Нехај \(\mathfrak{Q}\) значи квоциентно полье \(\mathfrak{I}\), а \(\mathfrak{Z}\) интегритетну област меджу \(\mathfrak{I}\) и \(\mathfrak{Q}\). Како выше, \(M\) именује се \textit{максимална} в \(\mathfrak{Z}\), ако из принадлежности \(a\cdot g(x)\), где \(a\ne0\) јест цело число и \(g(x)\in\mathfrak{Z}\), всегда следује принадлежност \(g(x)\). К \(\mathfrak{I}\) једнозначно ексистује максимална област \(M(\mathfrak{I})\), породжена чрез \(\mathfrak{I}\) в \(\mathfrak{Z}\), како пресек всих максималных областиј в \(\mathfrak{Z}\), кторе обхватајут \(\mathfrak{I}\); бо сама \(\mathfrak{Z}\) јест така, и с либокако многими је такоже пресек максималны и обхвата \(\mathfrak{I}\). \(M(\mathfrak{I})\) состојí из всих функциј \(g(x)\) из \(\mathfrak{Z}\), за кторе јест најменеј једно \(a\ne0\) с \(a\cdot g(x)\in\mathfrak{I}\).

Бо те \(g(x)\) творет максималну област \(N\) в \(\mathfrak{Z}\), ктора обхвата \(\mathfrak{I}\): из \(b\cdot h(x)\in\mathfrak{I}\), \(b\ne0\), \(h(x)\in\mathfrak{Z}\), следује \(ab\cdot h(x)\in\mathfrak{I}\) и \(ab\ne0\), зато \(h(x)\in N\). Всака максимална област \(M\) в \(\mathfrak{Z}\), ктора обхвата \(\mathfrak{I}\), обхвата \(N\); бо ако \(a\cdot g(x)\in\mathfrak{I}\), тогда \(a\cdot g(x)\in M\), понеже \(\mathfrak{I}\subset M\), и зато \(g(x)\in M\).

\medskip\noindent\textbf{2.}
Тепер нехај \(\mathfrak{I}\) буде предпоставјена како конечна интегритетна област с јединицоју, с \(f_1(x),\ldots,f_r(x)\) како интегритетноју базоју; то јест \(\mathfrak{I}\) состојí из всих целочислных полиномов од \(f_i(x)\). Далеј нехај \(\mathfrak{Z}\) состојí из всих целочислных полиномов или степеновых редов из \(x\), респективно из такых квоциентов, кторе не имајут инех именовательев кроме тых, кторе јавјајут се в конечно многих \(f_i(x)\); ты наравно в всаком степену.

Далеј предпоставльа се, же в најменеј једној \(f_i(x)\), обчно в всих, \(x\) стварно наступа; иначе \(\mathfrak{I}\) и \(\mathfrak{Z}\) просто состојат из всих целочислных полиномов ломаного числа \(1/e\), и ничто не треба доказати.

Ако идé о полиномы или рационалне функције, тогда квоциентно полье \(\mathfrak{Q}\) има трансценденцијны степен \(t\), \(1\le t\le r\), гледе на полье \(K\) рационалных чисел. Ако, например, \(f_1(x),\ldots,f_t(x)\) творет иредуцибилну систему\footnote{Под "иредуцибилноју системоју" разумеје се, како познано, система из алгебраично независимых функциј.}, тогда \(\mathfrak{Q}\) јест конечно алгебраично разширенье чисто трансценденцијного польа \(K(f_1(x),\ldots,f_t(x))\). Једновременно функционална матрица \(f_1(x),\ldots,f_r(x)\), и такоже она од \(f_1(x),\ldots,f_t(x)\), има точно ранг \(t\)\footnote{То додатно условје јест изполньено, ако идé, како в случају Хецке, о аналитичне функције, \(f_1(x),\ldots,f_t(x)\) сут аналитично независиме, а \(f_{t+1}(x),\ldots,f_r(x)\) алгебраично зависе од ньих. Але условје намерно формуловано јест формално --- ранг функционалној матрице, --- абы оно имало јасны смысл и при формално дефинованых степеновых редах. Тако все разсудживаньа оставајут чисто формално-алгебраичне, а теоремы о аналитичных функцијах потребне сут само за показанье условја в специалном случају.}.

Абы и в случају степеновых редов или јих квоциентов та структура \(\mathfrak{Q}\) остала, треба додати условје: ако \(t\) јест ранг функционалној матрице \(f_i(x)\) и напр. једновременно ранг функционалној матрице \(f_1(x),\ldots,f_t(x)\), тако же оне по §1, 1 творет иредуцибилну систему, тогда \(\mathfrak{Q}\) знову должно быти конечно алгебраично разширенье чисто трансценденцијного \(K(f_1(x),\ldots,f_t(x))\); то јест \(f_{t+1}(x),\ldots,f_r(x)\) должне сут быти алгебраичне гледе на \(K(f_1(x),\ldots,f_t(x))\)\footnote{То додатно условје изполньено јест в аналитичном случају Хецке: \(f_1(x),\ldots,f_t(x)\) сут аналитично независиме, а \(f_{t+1}(x),\ldots,f_r(x)\) алгебраично зависе од ньих.}. При том знову \(t\ge1\), бо в карактеристики нула не могу изчезати все \(\partial f_i/\partial x_k\).

Једновременно условје следује за всаку систему из \(t\) функциј интегритетној базы, кторых функционална матрица има точно ранг \(t\); бо те творет иредуцибилну систему, од ктореј остале алгебраично зависе.

\medskip\noindent\textbf{3.}
Нехај тепер \(\mathfrak{o}\) значи целочислну полиномну област в неизвестных \(u_1,\ldots,u_r\). Тогда \(\mathfrak{I}\) јест колцово-хомоморфны образ \(\mathfrak{o}\), како непосредно показује присвојенье \(u_i\mapsto f_i(x)\). Зато \(\mathfrak{I}\) јест колцово изоморфна к \(\mathfrak{o}/\mathfrak{a}\), кде \(\mathfrak{a}\) јест просты идеал, бо \(\mathfrak{I}\) јест колцо без нуловых делительев; \(\mathfrak{a}\) состојí из всих полиномов \(G(u)\), за кторе \(G(f)\) идентично изчезаје в \(x\).

Како уже споменуто в уводу, идé о теорем.

\medskip\noindent\textbf{Теорем.}
\textit{Под додатным условјем из 2. нехај \(\mathfrak{o}_p,\mathfrak{a}_p,\mathfrak{I}_p\) значет области, кторе из \(\mathfrak{o},\mathfrak{a},\mathfrak{I}\) настајут заменоју целочислных коефициентов јих остатковыми класами модуло \(p\). Тогда, с изкльученьем највыше конечно многих првочисел, меджу кторыми сут ты в именовательах \(f(x)\), абы \(\mathfrak{I}_p\) была дефинована, хомоморфизм од \(\mathfrak{o}_p\) к \(\mathfrak{I}_p\) посредованы јест чрез \(\mathfrak{a}_p\). Зато \(\mathfrak{I}_p\) јест колцово изоморфна к \(\mathfrak{o}_p/\mathfrak{a}_p\), и тем к \(\mathfrak{o}/(\mathfrak{a},p\mathfrak{o})\).}

Из теорема следује факт о само конечно многих првочислах в именовательах \(a\) области \(M(\mathfrak{I})\) чрез

\medskip\noindent\textbf{Последок.}
\textit{Ако \(p\) избрано јест тако, же \(\mathfrak{I}_p\) изоморфна јест к \(\mathfrak{o}_p/\mathfrak{a}_p\), тогда из \(p\cdot g(x)\in\mathfrak{I}\) и \(g(x)\in\mathfrak{Z}\) всегда следује \(g(x)\in\mathfrak{I}\).}

Другыми словами: \(\mathfrak{I}\) јест максимална в \(\mathfrak{Z}\) гледе всих првочисел кроме конечно многих изкльученых првочисел.

\(\mathfrak{o}_p/\mathfrak{a}_p\) јест изоморфно к \(\mathfrak{o}/(\mathfrak{a},p\mathfrak{o})\). Бо по дефиницији \(\mathfrak{o}_p=\mathfrak{o}/p\mathfrak{o}\), и соответне \(\mathfrak{a}_p=(\mathfrak{a},p\mathfrak{o})/p\mathfrak{o}\). Зато по првом изоморфизмном теорему \(\mathfrak{o}_p/\mathfrak{a}_p\), и по предпоставјеньу такоже \(\mathfrak{I}_p\), изоморфне сут к \(\mathfrak{o}/(\mathfrak{a},p\mathfrak{o})\).

То значи: из \(H(f_i(x))\equiv0\pmod p\) в \(\mathfrak{I}\) следује \(H(u)\equiv0\pmod{(\mathfrak{a},p\mathfrak{o})}\) в \(\mathfrak{o}\).

То можно преформовати так: из \(H(f(x))=p\cdot g(x)\), кде \(H(f(x))\in\mathfrak{I}\) и \(g(x)\in\mathfrak{Z}\), следује \(H(u)-pF(u)\equiv0\pmod{\mathfrak{a}}\), зато такоже \(H(f_i(x))=pF(f_i(x))\), и тем \(g(x)=F(f_i(x))\); зато \(g(x)\in\mathfrak{I}\).

Конечно многократно повторјенье дава: ако \(a\) не делимо јест ниједным изкльученым првочислом, тогда из \(a\cdot g(x)\in\mathfrak{I}\) и \(g(x)\in\mathfrak{Z}\) всегда следује \(g(x)\in\mathfrak{I}\). Последок доказаны.

\bigskip
\noindent\textbf{§ 3. Доказ теорема.}

\smallskip
\noindent\textbf{1. Помочна лема.}
\textit{Просты идеал \(\mathfrak{a}\), кторы посредује хомоморфизм меджу \(\mathfrak{I}\) и \(\mathfrak{o}\) (§2, 3), јест абсолутны просты идеал.}

Нехај \(\mathfrak{o}^*\) значи полиномну област в \(u_1,\ldots,u_r\) с либокаквыми, такоже ломаними, алгебраичными числами како коефициентами, а \(\mathfrak{a}^*\) идеал в \(\mathfrak{o}^*\), породжены чрез \(\mathfrak{a}\). Он состојí из всих линеарных комбинациј \(a_1G_1(u)+\cdots+a_sG_s(u)\), где \(a_i\) сут либокакве алгебраичне числа, а \(G_i(u)\) полиномы из \(\mathfrak{a}\). Треба доказати, же \(\mathfrak{a}^*\) јест просты идеал.

То буде доказано, ако показано буде, же чрез \(\mathfrak{a}^*\) посредује се хомоморфизм од \(\mathfrak{o}^*\) к \(\mathfrak{I}^*\), под \(\mathfrak{I}^*\) разумевши интегритетну област всих полиномов од \(f_i(x)\) с либокаквыми алгебраичными коефициентами; то јест, же \(H(u)\) принадлежи \(\mathfrak{a}^*\), ако \(H(u)\in\mathfrak{o}^*\) и \(H(f_i(x))=0\).

Нехај \(\omega_1,\ldots,\omega_s\) значи линеарно независиму базу конечного числового польа, породженого конечно многими коефициентами \(H(u)\). Тогда
\[
a\cdot H(u)=\omega_1H_1(u)+\cdots+\omega_sH_s(u),
\]
кде \(H_i(u)\in\mathfrak{o}\), а \(a\ne0\) јест цело рационално число. Из \(H(f)=0\) следује \(\omega_1H_1(f)+\cdots+\omega_sH_s(f)=0\), и зато \(H_1(f)=0,\ldots,H_s(f)=0\). Зато \(H_1(u),\ldots,H_s(u)\) принадлежет \(\mathfrak{a}\), и тем \(H(u)\in\mathfrak{a}^*\).

\medskip\noindent\textbf{2.}
Доказ теорема можно тепер дати в следујучној форме.

Ако \(p\) избрано јест тако, же
\begin{itemize}
\item[И.] \(\mathfrak{I}_p\) дефинована јест,
\item[ИИ.] \(\mathfrak{a}_p\) остава просты, именно абсолутны, идеал,
\item[ИИИ.] трансценденцијны степен \(\mathfrak{I}_p\) проти тому од \(\mathfrak{I}\) не понизил се,
\end{itemize}
тогда хомоморфизм од \(\mathfrak{o}_p\) к \(\mathfrak{I}_p\) посредованы јест чрез \(\mathfrak{a}_p\). Ты условја, с изкльученьем највыше конечно многих првочисел, изполньене сут за все првочисла.

Же изкльученых првочисел гледе И, ИИ, ИИИ јест само конечно мнохо, следује непосредно из предходного. И јест јасно: изкльученье сут само конечно многе првочисла в именовательах \(f_i(x)\). ИИ следује из лемы, же \(\mathfrak{a}\) јест абсолутны просты идеал и зацховава то својство модуло всих првочисел кроме највыше конечно многих. ИИИ показано јест в последку §1, 2, бо функционална матрица \(f_1(x),\ldots,f_t(x)\) по §2, 2 ма точно ранг \(t\).

Нехај тепер \(p\) избрано јест мимо тих изкльученых првочисел. Хомоморфизм од \(\mathfrak{o}_p\) к \(\mathfrak{I}_p\) посредованы јест идеалом \(\mathfrak{c}_p\), кторы јест просты, бо \(\mathfrak{I}_p\) такоже јест колцо без нуловых делительев, и јего трансценденцијны степен даны јест трансценденцијным степеньем \(\mathfrak{I}_p\), зато најменеј \(t\). Бо \(\mathfrak{I}_p\simeq\mathfrak{o}_p/\mathfrak{c}_p\), остатково полье \(\mathfrak{c}_p\) изоморфно јест к квоциентному польу \(\mathfrak{I}_p\).

Како \(\mathfrak{a}_p\) јест такоже просты идеал и многократник, то јест подидеал, од \(\mathfrak{c}_p\), тогда \(\mathfrak{a}_p\) и \(\mathfrak{c}_p\) соглашајут се точно тогда, когда јих трансценденцијне степены соглашајут се\footnote{Срав. например Б. Л. ван дер Wаерден, Зур Нуллстеллентхеорие дер Полыномидеале, Матх. Анн. 96 (1926), с. 183--208, §3.}. Трансценденцијны степен \(\mathfrak{a}_p\) како многократника од \(\mathfrak{c}_p\) јест такоже најменеј \(t\); класе \(u_1,\ldots,u_t\) модуло \(\mathfrak{a}_p\) творет иредуцибилну систему. Треба показати, же тој степен јест точно \(t\); с тым все буде доказано.

Туто в случају степеновых редов или јих квоциентов употребльа се додатно условје о структуре квоциентного польа \(\mathfrak{Q}\) области \(\mathfrak{I}\) (§2, 2). По тому просты идеал \(\mathfrak{a}\) в всаком случају имал трансценденцијны степен \(t\), и нумерација была избрана тако, же класе \(u_1,\ldots,u_t\) модуло \(\mathfrak{a}\) творили иредуцибилну систему, од ктореј класе \(u_{t+1},\ldots,u_r\) алгебраично зависели. В \(\mathfrak{a}\) были зато целочислне и примитивне полиномы
\[
G_1(u_{t+1};u_1,\ldots,u_t),\ldots,G_{r-t}(u_r;u_1,\ldots,u_t),\tag{1}
\]
кторе преходет в, због примитивности неизчезајуче, полиномы
\[
G^*_1(u_{t+1};u_1,\ldots,u_t),\ldots,G^*_{r-t}(u_r;u_1,\ldots,u_t),\tag{2}
\]
из \(\mathfrak{a}_p\). Тут але в \(G^*_j\) мора стварно наступати соответны \(u_{t+j}\), бо иначе, противно избору \(p\), бы настала алгебраична зависимост меджу класами \(u_1,\ldots,u_t\) модуло \(\mathfrak{a}_p\). Тым показано јест, же \(\mathfrak{a}_p\) има трансценденцијны степен \(t\), зато јест идентична с \(\mathfrak{c}_p\) и посредује хомоморфизм. Все доказано.

\bigskip
\noindent\textbf{§ 4. Разширенье на целе алгебраичне коефициенты.}

\smallskip
Како уже на концу увода примечено, все разсудживаньа с малыми модификацијами оставајут валидне, ако наместо целых рационалных чисел взяти целе числа конечного алгебраичного числового польа \(K\), а наместо првочисел \(p\) просте идеалы \(\mathfrak{p}\) того числового польа. При том §1 остава полно валидны, и такоже разсудживаньа, кторе ведут к формулацији теорема в §2; при доказу (§3, 2) и при последку в §2, 3 потребне сут неке додаткы.

\medskip\noindent\textbf{1.}
В §3, 2 треба к изкльученым простым идеалам додати јеште конечно многе, кторе вцходет в полиномы \(G_1(u),\ldots,G_{r-t}(u)\), §3, 2, (1), бо те тепер не можно предпоставльати примитивными. Тым гарантовано јест неизчезньенье полиномов \(G^*_1(u),\ldots,G^*_{r-t}(u)\), §3, 2, (2); то јест заоштренье условја ИИИ. Условје ИИ остава, бо теорем о абсолутној иредуцибилности не меньа се; такоже наравно остава И.

\medskip\noindent\textbf{2.}
Последок из §2, 3, кторы дава конечны резултат, треба соответне к тамој конечној форме формуловати тако:

\noindent\textbf{Последок.}
\textit{Ако цело число \(a\) из \(K\) не делимо јест ниједным из конечно многих изкльученых простых идеалов, тогда из \(a\cdot g(x)\in\mathfrak{I}\) и \(g(x)\in\mathfrak{Z}\) всегда следује \(g(x)\in\mathfrak{I}\).}

Ако \(\lambda_1,\lambda_2,\ldots\) сут целе числа из \(K\), избране тако, же всако делимо јест једным изкльученым простым идеалом, але не јего квадратом и не другыми, тогда за именователье \(a\) достачујут степенове продукты тих \(\lambda\).

\noindent\textit{Доказ.}
Нехај
\[
(a)=\mathfrak{p}_1^{\alpha_1}\cdots\mathfrak{p}_g^{\alpha_g}
\]
јест разложенье \((a)\) на просте идеалы и \(\mathfrak{p}_1,\ldots,\mathfrak{p}_g\) сут различне од изкльученых простых идеалов, зато \(\mathfrak{I}_{\mathfrak{p}_i}\) изоморфна јест к \(\mathfrak{o}_{\mathfrak{p}_i}/\mathfrak{a}_{\mathfrak{p}_i}\) за всакы \(\mathfrak{p}_i\). Од колца всих целых чисел из \(K\) преходимо к квоциентному колцу \(R_a\), дефинованому чрез \((a)\), додајучи в именователь все и само ты целе числа, кторе сут приме к \((a)\), то јест не делиме ниједным из идеалов \(\mathfrak{p}_1,\ldots,\mathfrak{p}_g\). В \(R_a\) идеалы \(\mathfrak{p}_1,\ldots,\mathfrak{p}_g\) оставајут просте, але стајут главне; все друге стајут једнотным идеалом; разложенье \((a)\) остава исто\footnote{Срав. например Х. Грелл, Зур Тхеорие дер Орднунген ин алгебраисцхен Захл- унд Функтионенкöрперн, Матх. Анн. 97 (1927), с. 524--558, §2, 1.}. Зато при преходу к \(R_a\) како коефициентној области последок и јего доказ оставајут точно в форме §2, 3.

То значи: ако \(a\cdot g(x)\in\mathfrak{I}\), тогда \(g(x)=F(f_i(x))\), кде полином \(F(u)\) има коефициенты из \(R_a\). Ако \(d\) јест сполечны именователь тих коефициентов, зато \(d\) јест примы к \(a\), тогда можно формуловати и тако: за \(a\cdot g(x)\in\mathfrak{I}\) лежи \(d\cdot g(x)\in\mathfrak{I}\) с \(d\) примым к \(a\). Зато \((d,a)\) јест једнотны идеал, и тем из \(a\cdot g(x)\in\mathfrak{I}\) следује \(g(x)\in\mathfrak{I}\). Прво тврдженье последка доказано.

За доказ другого тврдженьа нехај \(R^*\) јест квоциентно колцо в \(K\), дефиновано конечно многыми изкльучеными простыми идеалами, и \(\lambda_1,\ldots,\lambda_s\) нехај сут базне елементы тих простых идеалов в \(R^*\), кторе можно взяти како целе числа. В \(R^*\) всако \(a\) јест, до јединице из \(R^*\), степеновы продукт \(\lambda\). Врачајучи се к целым числам добивамо
\[
\gamma a=\beta\lambda_1^{\alpha_1}\cdots\lambda_s^{\alpha_s},
\]
кде \(\gamma\) и \(\beta\) не делиме сут ниједным изкльученым простым идеалом. Из \(a\cdot g(x)\in\mathfrak{I}\) зато следује, же \(\beta\lambda_1^{\alpha_1}\cdots\lambda_s^{\alpha_s}g(x)\), и по првом тврдженьу последка такоже \(\lambda_1^{\alpha_1}\cdots\lambda_s^{\alpha_s}g(x)\), лежи в \(\mathfrak{I}\).

\begin{flushright}
\textit{(Пријето 11 новембра 1931.)}
\end{flushright}

\bigskip
\noindent\rule{\textwidth}{0.4pt}

\medskip
\noindent\textbf{Абстракт} (превод из руского резјуме в оригиналу):

\smallskip
\textit{Всака интегритетна област \(\mathfrak{I}\) дефинује неку максималну област \(M(\mathfrak{I})\) истого типа, составјену из всих функциј \(g(x)\), за кторе јест најменеј једно \(a\ne0\) тако, же \(a\cdot g(x)\) садржано јест в \(\mathfrak{I}\).}

\textit{В тој праци изследује се впрашанье о "именовательах" \(a\) тој максималној области \(M(\mathfrak{I})\), под предпоставјеньем же област \(\mathfrak{I}\) была конечна. Доказано јест, же те \(a\) можно всегда избирати тако, абы в ньих скупај вцходило само конечно число првочисел како факторов.}

\textit{Једновременно в случају, кде \(g(x)\) сут степенове реды или квоциенты такых редов, предпоставльа се, же функције \(f_i(x)\), кторе дефинују област \(\mathfrak{I}\), могу быти меджусобно повезане само алгебраичными релацијами.}

\textit{Доказ опира се на свод проблему к некојим впрашаньам обчној теорије идеалов. Тен свод изполньа се разматраньем појавјеньа првочисла \(p\) в именовательу \(a\) како релације модуло \(p\) меджу функцијами \(f_i\).}

\end{document}
